\section{الفصل~\ref{sec:debugging}. تصحيح أخطاء الآلة}

ما يسمعه القطب، وانخفاض أبعاد النشاط، وعمل الواجهات القائمة على المصفوفات الصلبة، والتعلم المشترك للدماغ والمفكِّك، والنقاط البصرية الناتجة عن التحفيز، كلها مُقاسة في تجارب محكَّمة؛ وتقديرات تدفقات البيانات حساب في الفصل؛ أما الجهاز المزروع في البشر من \foreignlanguage{english}{Neuralink}، فبقدر ما أمكن التحقق، نشرت بياناته الشركة نفسها في الغالب.

{\footnotesize
\setlength{\tabcolsep}{3pt}\renewcommand{\arraystretch}{1.15}
\begin{longtable}{>{\raggedright}p{0.5cm}>{\raggedright}p{4.0cm}>{\raggedright}p{5.6cm}>{\raggedright}p{2.3cm}>{\raggedright\arraybackslash}p{1.7cm}}
\toprule
م & القضية & التجربة وما قيس & المصدر & الحالة \\
\midrule\endfirsthead
\toprule
م & القضية & التجربة وما قيس & المصدر & الحالة \\
\midrule\endhead
1 & يسمع القطب نبضات العصبونات في نصف قطر من رتبة مئة ميكرومتر، عادةً من عصبون واحد إلى بضعة عصبونات؛ وتُميَّز بشكلها & جرذ، الحقل \foreignlanguage{english}{CA1}، تسجيل داخل خلوي وخارج خلوي متزامن: شكل النبضة خارج الخلية يكرر عرض النبضة داخلها وسعتها. تقدير: كان يمكن للقطب الرباعي أن يميز $60\text{--}100$ عصبون، لكن ما يُستخرج فعلًا نحو ستة. & \cite{Henze2000extra} & مُقاس \\
2 & في الدماغ $8.6\cdot10^{10}$ عصبون؛ والمسبار يسمع نحو جزء من مئة مليون & عدّ النوى في معلّق من النسيج المذاب: $86\cdot10^9$ عصبون في دماغ الإنسان. والنسبة حساب في الفصل. & \cite{Azevedo2009} & مُقاس \\
3 & نشاط مئات العصبونات في القشرة الحركية ينحصر في عشرة أو عشرين متغيرًا مشتركًا & مراجعة لتسجيلات في القشرة الحركية عند القردة: نشاط المجموعة توصفه متغيرات كامنة قليلة مشتركة بين عصبونات كثيرة. & \cite{Gallego2017} & مُقاس \\
4 & ضجيج الجيران متشارك، ومئة عصبون تحمل تقريبًا ما تحمله ألف (القضية~\ref{prop:pooling}) & القشرة البصرية عند القرد: معامل ارتباط ضجيج العصبونات المتجاورة نحو $0.1$، والمكسب من المتوسط يتشبع عند مئة عصبون. ونظرية الارتباطات المقيِّدة للمعلومات. & \cite{Zohary1994, MorenoBote2014} & مُقاس \\
5 & التدفق الخام $2\cdot10^8$ بت/ثانية مقابل $8\cdot10^4$ بت/ثانية في النبضات~\eqref{eq:probedata} & حساب في الفصل من سعة تدفق النبضات~\eqref{eq:spikeentropy}. ومعاملات الرقمنة واقعية: رقاقة \foreignlanguage{english}{Neuralink} بتردد $19.3\,\text{kHz}$ و$10$ بتات لكل قناة. & اشتقاق في الفصل؛ \cite{Rieke1997, Musk2019} & نظرية \\
6 & المنظومة الأولى لـ\foreignlanguage{english}{Neuralink}: الرقاقات تضخّم وترقمن، والتدفق الخام يخرج عبر كابل، وبرنامج في الخارج يستخرج النبضات؛ أما الاستخراج على الرقاقة والغلاف المغلق والاتصال اللاسلكي والطاقة بلا أسلاك فللجهاز المخصص للبشر، بحسب ما أعلنته الشركة & مقالة الشركة: التدفق الخام الكامل عبر كابل \foreignlanguage{english}{USB-C}، وبرنامج يستخرج النبضات في الزمن الحقيقي خارج الرقاقة؛ والضغط على متن الجهاز والطاقة والإرسال بلا أسلاك ذُكرت خطةً للأجهزة السريرية. وعن الجهاز المزروع في البشر لا يوجد إلا إعلانات الشركة. وضرورة استخراج النبضات على الرقاقة استنتاج الفصل من~\eqref{eq:probedata}. & \cite{Musk2019}؛ تقارير \foreignlanguage{english}{Neuralink}، لا مقالة محكَّمة & مُقاس (مقالة الشركة)؛ وللبشر تقرير الشركة \\
7 & حتى $3072$ قطبًا على $96$ خيطًا، $32$ لكل خيط؛ يغرس الخيوطَ روبوتٌ يتجنب الأوعية & مقالة الشركة: $3072$ قطبًا على $96$ خيطًا، $32$ لكل خيط؛ يغرس الروبوت $6$ خيوط ($192$ قطبًا) في الدقيقة؛ وحتى $70\,\%$ من الأقطاب لدى جرذان بمصفوفة مزروعة لمدة طويلة تسمع نبضات. & \cite{Musk2019} & مُقاس (مقالة الشركة) \\
8 & عند البشر نحو ألف قناة على $64$ خيطًا؛ بعض الخيوط انزاح؛ والمؤشر نحو $10$ بت/ثانية & إعلانات الشركة فقط. لم يعثر البحث في \foreignlanguage{english}{PubMed} على مقالة محكَّمة بنتائج على البشر؛ وهذا يتفق مع التحفظ في نص الفصل. & تقارير \foreignlanguage{english}{Neuralink}؛ لا مقالة محكَّمة & مُقاس (تقرير الشركة) \\
9 & واجهة على مصفوفة من مئة قطب تتحكم في المؤشر & إنسان مشلول، $96$ قطبًا في القشرة الحركية الأولية: نية تحريك الذراع تغيّر النبضات بعد ثلاث سنوات من الإصابة؛ ويعطي المفكِّك ”مؤشرًا عصبيًا“ (البريد، التلفاز)، والتحكم في يد اصطناعية وفي ذراع روبوتية. & \cite{Hochberg2006} & مُقاس \\
10 & الكتابة ”باليد المتخيَّلة“: نحو $90$ حرفًا في الدقيقة، أقل من $8$ بت/ثانية & إنسان مشلول اليد، محاولات الكتابة باليد، ومفكِّك تكراري: $90$ حرفًا في الدقيقة بدقة $94.1\,\%$ في الزمن الحقيقي، وأكثر من $99\,\%$ بعد التصحيح التلقائي. والتقدير بالبتات حساب في الفصل. & \cite{Willett2021} & مُقاس \\
11 & محاولات الكلام تتحول إلى نص بسرعة نحو $60$ كلمة في الدقيقة & إنسان فقد الكلام الواضح، مصفوفات في القشرة: $62$ كلمة في الدقيقة؛ و$9.1\,\%$ أخطاء في الكلمات بمعجم من $50$ كلمة، و$23.8\,\%$ بمعجم من $125\,000$. & \cite{Willett2023} & مُقاس \\
12 & تستخرج الواجهة جزءًا صغيرًا من السعة: عشرات البتات في الثانية للعصبون الواحد، وآلاف للمئة (القضية~\ref{prop:probe}) & الحد الأعلى~\eqref{eq:spikeentropy} نظرية؛ والمقارنة بالسطرين~8 و10 استنتاج في الفصل. أما أن عنق الزجاجة هو انخفاض الأبعاد والجهل بالشفرة لا السلك، فلم يُختبر بتجربة مباشرة. & \cite{Rieke1997} & نظرية \\
13 & الدماغ والمفكِّك يتعلم كل منهما من الآخر & قرود تتحكم في المؤشر؛ المفكِّك يتكيف في حلقة مغلقة، والعصبونات يُعاد تعلّمها في الوقت نفسه: تزداد الدقة، وتبقى المهارة ويقل تأثرها بحركات اليد نفسها. ونصيحة الفصل بين سرعتي التعلم قاعدة هندسية، ولا تجربة مستقلة لها في الفصل. & \cite{Orsborn2014coad} & مُقاس \\
14 & التيار عبر القطب يثير العصبونات والأسلاك المحيطة دون تمييز للنوع & قشرة الفأر والقط، تسجيل الكالسيوم بالفوتونين: حتى التيار الضعيف يثير عصبونات متفرقة منتشرة على مسافة تصل إلى المليمترات، عبر إثارة المحاوير مباشرة في حجم قطره عشرات الميكرومترات. أي أن الإثارة غير انتقائية، لكنها ليست متصلة. & \cite{Histed2009stim} & مُقاس \\
15 & تحفيز القشرة البصرية يضيء نقاطًا؛ وحتى $16$ نقطة معًا تتيح التعرف على بعض الحروف وحدود الأشياء & إنسان كفيف، $96$ قطبًا في القشرة البصرية لمدة $6$ أشهر: عتبة القطب الواحد $66.8\pm36.5\,\text{\textmu A}$؛ والتحفيز المتزامن للمجموعات (حتى $16$ قطبًا، وهو حد المحفِّز) يخفض العتبة ويعطي أنماطًا قابلة للتمييز، ويُتعرَّف على بعض الحروف وحدود الأشياء. & \cite{Fernandez2021} & مُقاس \\
16 & الاتجاه الثاني لـ\foreignlanguage{english}{Neuralink} جهاز يحقن إشارات في القشرة البصرية & إعلانات الشركة عن التطوير فقط؛ لم يُعثر على بيانات محكَّمة عن عمله. & تقارير \foreignlanguage{english}{Neuralink} & فرضية \\
17 & تراكيز نواقل البث العصبية يمكن قياسها بمستشعر كيميائي في النسيج & شرائح دماغ الجرذ، قياس الفولتية الدوري السريع على ليف كربوني: قيس إطلاق السيروتونين وإعادة امتصاصه؛ والمادة تخرج خارج حدود المشبك. & \cite{Bunin1998} & مُقاس \\
\bottomrule
\end{longtable}}
